% answers_short.tex —— 2016 年·江苏卷 · 简短答案速查
% 由 tools/make_short_answers.py 自 answers.tex 的 \gksolhead/\gkanswer 提取，
% 并与 2016/2016_数学真题_江苏_答案要点.md 逐题比对后生成，勿手改（改源 MD 后重跑）。

\begin{gknotice}
\textbf{使用说明}：本节只给最终答案便于快速核对，完整推导见解析版。答案为编写者独立推导，\textbf{非官方标准答案}。
\end{gknotice}

\gkpart{填空题}
\begin{itemize}[leftmargin=2.6em]
  \item[\textbf{1.}] \(\{-1,2\}\)
  \item[\textbf{2.}] \(5\)
  \item[\textbf{3.}] \(2\sqrt{10}\)
  \item[\textbf{4.}] \(0.1\)
  \item[\textbf{5.}] \([-3,1]\)
  \item[\textbf{6.}] \(9\)
  \item[\textbf{7.}] \(\dfrac{5}{6}\)
  \item[\textbf{8.}] \(20\)
  \item[\textbf{9.}] \(7\)
  \item[\textbf{10.}] \(\dfrac{\sqrt{6}}{3}\)
  \item[\textbf{11.}] \(-\dfrac{2}{5}\)
  \item[\textbf{12.}] \(\left[\dfrac{4}{5},13\right]\)
  \item[\textbf{13.}] \(\dfrac{7}{8}\)
  \item[\textbf{14.}] \(8\)
\end{itemize}

\gkpart{解答题}
\begin{itemize}[leftmargin=2.6em]
  \item[\textbf{15.}] (1) \(5\sqrt{2}\)；(2) \(\dfrac{7\sqrt{2}-\sqrt{6}}{20}\)
  \item[\textbf{16.}] (1) 见解析版；(2) 见解析版
  \item[\textbf{17.}] (1) \(312\ \mathrm{m}^{3}\)；(2) \(PO_{1}=2\sqrt{3}\ \mathrm{m}\)（此时仓库容积为 \(416\sqrt{3}\ \mathrm{m}^{3}\)）
  \item[\textbf{18.}] (1) \((x-6)^{2}+(y-1)^{2}=1\)；(2) \(y=2x+5\) 或 \(y=2x-15\)；(3) \(t\in[2-2\sqrt{21},\ 2+2\sqrt{21}]\)
  \item[\textbf{19.}] (1) ① \(x=0\)；② \(m\) 的最大值为 \(4\)；(2) \(ab=1\)
  \item[\textbf{20.}] (1) \(a_{n}=3^{\,n-1}\)；(2) 见解析版；(3) 见解析版
\end{itemize}

\gkpart{附加题}
\begin{itemize}[leftmargin=2.6em]
  \item[\textbf{21.}] 【B】$AB=\begin{bmatrix}1&\dfrac{5}{4}\\0&-1\end{bmatrix}$；【C】\(\dfrac{16}{7}\)；【A】【D】为证明题，见解析版。
  \item[\textbf{22.}] (1) \(y^{2}=8x\)；(2) ① 证明见解析版；② \(0<p<\dfrac{4}{3}\)
  \item[\textbf{23.}] (1) \(0\)；(2) 见解析版
\end{itemize}
